\section{The Testbed}
\label{sec:testbed}

The model is fully discrete, so it is fully simulable. We built a testbed, an open-source TypeScript library, the \texttt{@cluesurf/vibe} package~\cite{pollard2026vibetest}, available under the MIT License at \url{https://github.com/cluesurf/vibe}, that generates a discrete substrate, runs a local rule over it in discrete beats, and measures the emergent physics. Everything is finite and seeded, so every number reported here is a deterministic function of a stated seed and reproduces exactly.

\subsection{Design}

The library is organized as a small set of layers. \emph{Substrates} generate the mesh: Poisson-sprinkled Minkowski and curved spacetime, regular lattices, hyperbolic tilings, hyperbolic random graphs, and classical sequential growth. \emph{Rules} update it: synchronous, asynchronous, reversible, rewriting, and gauge updates. \emph{Operators} act on it: the graph Laplacian, the Kahler-Dirac operator, the evolution Hamiltonian, the gauge-covariant Dirac operator, and the overlap operator. \emph{Measures} read off physics: dimension, distance, curvature, manifold-likeness, Lorentz isotropy, navigability, the CHSH correlation, locality, integration, Wilson loops, and the chiral condensate. \emph{Dynamics} provide the samplers: the Benincasa-Dowker action, causal-set Monte Carlo, parallel tempering, exact enumeration, a correct uniform-measure sampler, and the Wilson heat bath.

Real numbers appear only as measured outputs, such as sprinkling coordinates and eigenvalues, never as the substrate. This keeps the discreteness commitment honest at the level of the code.

\subsection{Trustworthiness}

A model is only as trustworthy as its checks. The testbed carries a suite of known-answer tests that must pass before any result is believed: a Poisson sprinkling of a $d$-dimensional region must recover dimension $d$, a regular lattice must be more anisotropic than a sprinkling, a CHSH experiment with independent settings must respect the classical bound, the overlap operator must have one species with exact chiral symmetry, the uniform-measure sampler must reproduce exact enumeration, and so on. Several early results were wrong in plausible-looking ways and were caught only by these checks, including a miscalibrated dimension estimator, a sprinkler that was not uniform by volume, an anisotropy statistic that cancelled for symmetric lattices, and a Monte Carlo move that did not sample the uniform measure. Each was fixed and recorded. The results below are reported only for measures that pass their known-answer tests.

\subsection{The nine questions}

We pose the model's load-bearing questions as P1 through P9, each a runnable experiment:
\begin{itemize}
\item \textbf{P1.} Can the dynamics have a bounded-below and local Hamiltonian (a stable ground state and a local time)?
\item \textbf{P2.} Is there a dynamics under which manifold-like spacetime dominates?
\item \textbf{P3.} Can one substrate be at once Lorentz-safe and navigable?
\item \textbf{P4.} Is matter, including spin, a pattern of the mesh?
\item \textbf{P5.} Is the geometry recovered from a discrete order unique?
\item \textbf{P6.} Is the 2D sum over histories computable and 2-dimensional?
\item \textbf{P7.} Can quantum statistics come from the deterministic mesh?
\item \textbf{P8.} Can a gauge force and a charged fermion be built?
\item \textbf{P9.} What is the relation between structure and experience?
\end{itemize}
The next sections answer these in the order of the emergence ladder, grouping them by the structure they build.
